% ============================================================
% 片段：s00  题名页、摘要、译者说明、目录
% 内容：对应原文 main.tex 的题名页/摘要部分（PDF 第 1--3 页）
% 译者：lead   审校状态：初稿
% ============================================================

\begin{titlepage}
  \centering
  \vspace*{1.2cm}

  {\LARGE\bfseries 广义对称性及其应用讲义\par}
  \vspace{0.6cm}
  {\large Lecture Notes on Generalized Symmetries and Applications\par}
  \vspace{0.4cm}
  {\normalsize 中文学术译本（附审校译注与本节重点）\par}

  \vspace{1.6cm}
  {\large Ran Luo$^{1}$, \ Qing-Rui Wang$^{2}$, \ Yi-Nan Wang$^{1,3}$\par}

  \vspace{1.0cm}
  {\small
    $^1$ 北京大学物理学院，北京 100871\par
    \vspace{0.3cm}
    $^2$ 清华大学丘成桐数学科学中心，北京 100084\par
    \vspace{0.3cm}
    $^3$ 北京大学高能物理中心，北京 100871\par
  }

  \vspace{1.4cm}
  {\small arXiv:2307.09215v3 [hep-th] \quad 2024 年 8 月 23 日（v3，83 页）\par}
  \vspace{0.25cm}
  {\small 原载 \emph{Physics Reports}，doi:10.1016/j.physrep.2024.02.002\par}

  \vfill
  {\footnotesize
  \begin{minipage}{0.88\linewidth}
  \textbf{版权与分享声明。}原文 \copyright\ 2024 Ran Luo, Qing-Rui Wang, Yi-Nan Wang，
  经 arXiv 以“非独占分发许可”（arXiv.org perpetual, non-exclusive license）发布；该许可只授权
  arXiv 分发原文，并不授予第三方翻译、改编或再发布的权利，原文著作权仍归作者所有。
  本译本是原文的演绎作品：翻译、图表重制、译注与排版由 DeepSeek Harness 智能体集群完成；
  译注与“本节重点”系译者独立添加的评论，不代表原作者观点，译文讹误由译者负责。
  本译本可原样用于学术研究与教学，不得商业使用，不得未经许可改编、摘编或再发布（尤其不得投稿、出版、
  收录进教材）；此类使用须先取得原作者与本译者的书面许可。译注部分可单独按 CC BY-NC-SA 4.0 使用，
  但整份文档的许可上限由原文决定。若原作者或 arXiv 提出异议，将立即停止分发。
  \end{minipage}\par}
  \vspace{0.6cm}
\end{titlepage}

% ============================ 摘要 ============================
\thispagestyle{plain}
\begin{center}
  {\bfseries\large 摘\quad 要}
\end{center}
\vspace{0.3em}

在本讲义中，我们从高能物理与凝聚态物理两个视角，对广义对称性这一迅速发展的概念作基础性介绍。
我们特别强调（可逆的）高形式对称性与高群对称性。在物理应用方面，我们讨论弦论中量子场论的几何工程，
以及凝聚态物理中的对称保护拓扑（SPT）相。

本讲义基于王一男与王清睿于 2023 年 2 月在北京大学物理学院联合讲授的广义对称性短期课程
（\url{https://indico.ihep.ac.cn/event/18796/}）。

\vspace{1.2em}
\noindent\rule{\linewidth}{0.4pt}
\vspace{0.4em}

\noindent{\bfseries 译者说明}

\begin{itemize}[leftmargin=1.4em,itemsep=2pt,topsep=2pt]
  \item \textbf{底本}：译文以 arXiv:2307.09215v3（2024 年 8 月 23 日，83 页）为准，与
        \emph{Physics Reports} 正式发表版（doi:10.1016/j.physrep.2024.02.002）内容一致。
  \item \textbf{公式}：全部公式由原文 LaTeX 源码逐字搬运，公式环境与先后次序不变，
        故公式编号（按节编号，如 (2.26)、(5.12)）与原文逐一对应；
        具体设置见 \texttt{main.tex} 导言区的 \texttt{numberwithin} 一行。
  \item \textbf{图表}：Figure 1--7 由原文 PDF 以 400 dpi 精确裁切（\texttt{figures/}），图注中译并标注
        原文图号与 PDF 页码；Table 1 以及原文中不编号的第二张表用原生 \LaTeX{} 表格重排。
        正文中大量嵌入公式或句中的示意图（原文以 TikZ/tikzcd 绘制、不编号）保留其 TikZ 矢量源码，
        以保证与正文基线对齐且可编辑。
  \item \textbf{译注}：橙色“译注”方框为译者所加，全文连续编号，内容分三类：推导易错预警、
        作者省略之处（“it is obvious / one can show”）的点拨、原文疑误与排版问题说明。
        不需要译注时，把 \texttt{main.tex} 中的 \verb|\showyizhutrue| 改为 \verb|\showyizhufalse| 即可整体关闭。
  \item \textbf{本节重点}：每节开头的蓝色方框为译者所加的导读，不是原文内容。
  \item \textbf{编译}：需用 \textbf{XeLaTeX} 引擎（中文由 \texttt{ctex} 处理），连跑两遍：
        \texttt{xelatex main.tex}。在无 TeX~Live 的环境中也可用
        \texttt{tectonic -X compile main.tex}（XeTeX 引擎）。
  \item \textbf{文献}：参考文献沿用原文编号与英文著录（学术惯例），由原文 \texttt{main.bbl} 直接读入。
\end{itemize}

\vspace{0.6em}
\noindent{\small 译者（DeepSeek Harness 智能体集群）\quad 2026 年}

\clearpage
\tableofcontents
\clearpage
